\chapter{Consideraciones Finales}\label{cap:Considerações}

% ============================================================
%  BLOQUE C — IMPLICACIONES DINÁMICAS
% ============================================================

\section{Implicaciones dinámicas}

\textit{Primera.} La intensidad del vórtice central ($\zeta_{850}^{\min}$) predice adecuadamente la organización espacial del viento únicamente en ciclones ordinarios. En sistemas severos, el campo cinemático superficial queda controlado por procesos de mesoescala ---fractura frontal, CCB, \textit{sting jets}--- que no escalan linealmente con la vorticidad. Parametrizar el viento destructivo a partir de métricas de vorticidad mínima subestimará sistemáticamente la magnitud y localizará incorrectamente los máximos durante la fase crítica de madurez.

\textit{Segunda.} Los ciclones que generan los vientos más destructivos no producen las precipitaciones más extremas en su fase madura. La relación vorticidad--precipitación cambia de signo en SBR y LPB durante la madurez porque la profundización del vórtice acelera la intrusión seca que suprime la convección central. Los máximos pluviométricos de los sistemas p90 se alcanzan durante la intensificación; en la madurez, la amenaza hidrológica se desplaza hacia la banda periférica enroscada en cuadrantes opuestos a los del viento máximo.

\textit{Tercera.} El desacople geométrico entre las amenazas de viento y precipitación no es la yuxtaposición accidental de dos patrones independientes. El MEOF1 de madurez del subconjunto p90 captura la anticovarianza entre cuadrantes como un único modo multivariado dominante, confirmando que la banda enroscada sureste y el núcleo de viento noroeste covarían sistemáticamente dentro de los mismos eventos. La separación de cuadrantes es la firma diagnóstica estructuralmente coherente y predecible del ciclón ocluido intenso del Atlántico Sur, no un promedio de configuraciones dispares.

\textit{Cuarta.} Las diferencias interregionales son físicamente sustantivas. SBR y LPB exhiben el desacople más pronunciado, sostenido por los flujos de calor latente oceánico y el aporte de humedad del SALLJ. ARG presenta oclusiones más graduales y un acoplamiento vorticidad--precipitación más persistente, reflejo de una dinámica baroclínica dominante menos modulada por procesos diabáticos locales. Estas diferencias son la expresión de regímenes dinámicos diferenciados, no artefactos estadísticos.

\textit{Quinta.} Hasta donde se ha podido verificar, este trabajo constituye el primer análisis que combina PDFe espacial, EOF univariado y MEOF para caracterizar conjuntamente el desacople viento--precipitación por fase del ciclo de vida en el Hemisferio Sur. Un estudio independiente y reciente, \citet{ribeiro2025impacts}, documenta mediante composites simples (sin EOF ni MEOF) una separación análoga entre los cuadrantes de viento (noreste) y precipitación (centro--sureste) durante la oclusión de ciclones en cuatro cuencas del Hemisferio Sur. Esta coincidencia, obtenida con una metodología independiente, corrobora la plausibilidad física del desacople geométrico aquí identificado y subraya, a la vez, la necesidad del enfoque multivariado más fino desarrollado en esta tesis para cuantificarlo por región, fase evolutiva e intensidad dinámica.

% ============================================================
%  BLOQUE D — LIMITACIONES
% ============================================================

\section{Limitaciones}

Los resultados están sujetos a la subestimación de ERA5 en extremos ($-10.2\%$ en viento y $-22.6\%$ en precipitación para los ciclones más intensos), por lo que las magnitudes absolutas del subconjunto p90 son estimaciones conservadoras. Las estructuras de mesoescala responsables de los máximos de viento en ciclones severos ---CCB, \textit{sting jets}, descargas descendentes--- operan a escalas menores que la resolución de ERA5 ($\sim$31~km), limitando la resolución de su geometría interna en los prototipos. La estandarización del dominio lagrangiano a $\pm 9{,}5^{\circ}$ puede truncar la banda de precipitación enroscada en los ciclones de mayor escala horizontal, subestimando la extensión real del campo pluviométrico periférico. El tamaño muestral reducido del subconjunto p90 (60--65 ciclones por región) introduce mayor incertidumbre en las densidades periféricas de las PDFe, aunque la significancia Bootstrap-$t$ al 95\% garantiza la robustez de las firmas espaciales centrales. El análisis se circunscribe a los campos de superficie; la extensión a la columna tridimensional permitiría vincular las estructuras identificadas con sus mecanismos generadores.

% ============================================================
%  BLOQUE E — PERSPECTIVAS FUTURAS
% ============================================================

\section{Perspectivas futuras}

\begin{enumerate}
    \item \textit{Validación con observaciones independientes.} La geometría de los prototipos debe contrastarse con observaciones de alta resolución (boyas, estaciones costeras, altimetría satelital) para cuantificar los sesgos de ERA5 y ajustar los valores modales de las PDFe en el subconjunto p90, en particular la extensión real de la banda enroscada de precipitación.

    \item \textit{Extensión tridimensional.} Incorporar los campos de viento en 850~hPa, temperatura potencial equivalente y vorticidad potencial al análisis MEOF permitirá vincular las firmas superficiales con la evolución tridimensional de las corrientes transportadoras (WCB, CCB) y la estructura del \textit{sting jet} en columna.

    \item \textit{Análisis de eventos compuestos.} Los prototipos proveen el marco de referencia para definir umbrales combinados sobre las PDFe de viento y precipitación, cuantificando la probabilidad de impactos compuestos en las costas de Brasil, Uruguay y Argentina con diferenciación explícita entre el peligro eólico (noroeste) y el hidrológico (sureste--sur) en función de la posición del ciclón respecto a la costa.

    \item \textit{Proyecciones climáticas.} Evaluar en simulaciones de alta resolución con escenarios SSP si el calentamiento global modifica la separación de cuadrantes documentada ---en particular la posición y extensión del núcleo de viento noroeste y la coherencia de la banda enroscada--- con implicaciones directas para la gestión del riesgo costero en el Atlántico Sudoccidental.
\end{enumerate}
%\chapter{Discussão}\label{cap:discussão}

Embora este estudo....
